\section{Policy Implications: The Relay Network as I2.0 Layer Zero}
\label{sec:policy}

The economic cases in the preceding sections are unlocked by two distinct investments:
the I2.0 managed lane program (\$253 billion portfolio, decade-scale construction timeline)
and the relay station network (\$40 million for NY-LA, \$2 billion nationally, buildable
in three to five years). These investments are complements, not substitutes. But they are
not symmetric in urgency or in what they deliver independently.

\subsection{What the Relay Network Delivers Without Managed Lanes}

The simulation tested four scenarios, of which two are relevant to the policy sequence
question: relay on current General Purpose lanes, and relay on I2.0 managed lanes.

\begin{table}[h]
\centering
\caption{Relay performance with and without I2.0 managed lanes}
\label{tab:relay_scenarios}
\begin{tabular}{lrrrr}
\toprule
Scenario & p50 (h) & p75 (h) & p95 (h) & 48h SLA \\
\midrule
Relay / GP lanes (buildable now)   & 40.2 & 45.1 & 58.3 & $\approx$88\% \\
Relay / I2.0 managed lanes         & 35.6 & 39.8 & 45.4 & 98.8\% \\
\bottomrule
\multicolumn{5}{l}{\footnotesize Source: \citet{ROUTE_C1}. GP = General Purpose (current infrastructure).} \\
\end{tabular}
\end{table}

The relay network on current infrastructure delivers an 88\% probability of meeting
a 48-hour service level agreement. This is not sufficient for most commercial supply chains
that require a 48-hour guarantee --- a carrier who promises 48-hour delivery and misses
it 12\% of the time cannot sustain that promise. But it is dramatically better than the
current solo/GP baseline (where 48-hour delivery is essentially never achievable), and
it is good enough for less time-sensitive shipments that prefer 48-hour truck over 87-hour
truck regardless of occasional misses.

The practical implication: relay stations can be built today, on existing interstate
infrastructure, and begin capturing a portion of the economic value described in this paper.
They are not contingent on the managed lane program being complete. The managed lane program,
when built, improves the relay SLA from 88\% to 98.8\% --- the difference between
a ``best effort'' service and a commercial commitment window.

This sequencing matters for policy. The federal government does not need to wait for the
\$253 billion I2.0 program to be authorized and funded before beginning relay network
development. The relay network is a distinct, smaller program that can proceed on a
faster track.

\subsection{Policy Recommendation 1: Designate National Freight Relay Zones}

The Federal Highway Administration has authority under the National Highway Freight
Program (established under the FAST Act, 2015, and continued under IIJA, 2021) to
designate freight-critical facilities and fund infrastructure improvements at those
facilities \citep{FHWA_freight2023}. We recommend that FHWA exercise this authority
to designate a network of National Freight Relay Zones at T1 hub locations along
the primary Interstate 2.0 corridors.

The designation would accomplish three things. First, it would signal to the private
sector that the relay hub locations are stable and investable, enabling carriers to
begin planning relay operations without fear that the designated locations will change.
Second, it would unlock National Highway Performance Program funding for station
infrastructure --- access roads, utilities, truck pad construction, and basic facility
development --- that reduces the private capital requirement for carriers. Third, it
would establish a planning baseline for state DOTs to incorporate relay stations into
their highway interchange improvement programs.

The cost to the federal government of designation is zero. The cost of NHPP-funded
station infrastructure at nine NY-LA locations is \$40 million --- already within the
NHPP's annual allocation for freight-specific projects in the affected states.

\subsection{Policy Recommendation 2: HOS Regulatory Clarification for Relay Operations}

Current Hours of Service regulations were designed for solo and team driving.
The relay model creates a new compliance question: when a driver hands off the truck
at a relay station and is transported back to her home base as a non-driving passenger,
does that count as on-duty time?

Under current FMCSA regulations, a driver who is in a commercial vehicle but not driving
is considered to be ``on duty'' unless she is in the sleeper berth. A relay driver
returning home as a passenger in a repositioning truck would not be in a sleeper berth
and would technically accumulate on-duty hours during the return trip --- potentially
approaching HOS limits before her next driving shift \citep{FHWA_freight2023}.

This is a compliance artifact, not a safety concern. A relay driver who has worked a
seven-hour leg and then spent five hours as a passenger in a comfortable cab seat is
not fatigued in any meaningful sense. FMCSA should issue a clarifying rule or guidance
document that treats relay driver repositioning travel as off-duty time, provided the
driver has no active driving or supervisory responsibilities during the return trip.
This regulatory clarification costs nothing and removes a significant barrier to relay
operation adoption by risk-averse carriers.

The 30-minute break requirement is a related issue. Under current rules, a driver must
take a 30-minute break after eight hours of cumulative on-duty time (not just driving time).
A relay driver who works a seven-hour driving leg has not triggered this requirement, but
if any pre-drive inspection, paperwork, or wait time at the relay station is included,
the eight-hour clock may be close. FMCSA should clarify that the relay station changeover
process does not trigger the 30-minute break requirement when the driving leg is under
eight hours and the driver goes off-duty at the station.

\subsection{Policy Recommendation 3: The Freight Network Operator Model}

The relay network's full economic value requires coordination across carriers. The driver
who boards the truck in Cleveland does not know or care which carrier owns the trailer.
She is a relay driver; her job is to move the truck from Cleveland to Chicago. The trailer
might belong to Werner, J.B. Hunt, or a California agricultural cooperative's private fleet.
The relay system works only if the handoff at each station is seamless regardless of
carrier identity.

This is a coordination problem that individual carriers cannot solve unilaterally. No single
carrier has enough volume to fill eight relay stations continuously; the economics require
sharing the station infrastructure across carriers. The parallel to the airline industry is
instructive: no individual airline could maintain a full hub operation at every airport.
The shared infrastructure model --- airports owned and operated independently of airlines,
with gate access sold on a per-use basis --- solved this coordination problem.

We recommend that FHWA, in partnership with the Surface Transportation Board and FMCSA,
develop a Freight Network Operator framework: a regulatory structure that authorizes neutral
relay station operators to coordinate driver handoffs across carriers, maintain shared
scheduling systems, and manage station infrastructure under common operational standards.
The precedent exists in rail: trackage rights agreements allow freight railroads to use
each other's infrastructure under regulatory supervision. The relay network needs a
trucking equivalent.

A freight network operator would be neither a carrier (it does not own freight) nor a
broker (it does not arrange shipments). It is an infrastructure operator --- analogous
to an airport authority or a port authority --- that provides the physical and
organizational platform on which carriers operate. The legal framework for this entity
does not currently exist in federal trucking regulation; creating it requires a rulemaking
or statutory clarification.

\subsection{The Cost-Benefit Summary}

The policy package described above --- relay zone designation, HOS clarification, and
freight network operator framework --- has a total federal direct cost of approximately
\$40--50 million for the NY-LA pilot, or \$2 billion for a national network. The annual
economic benefit from the five economic cases analyzed in this paper totals \$15--25 billion
per year at full adoption.

\begin{equation}
\text{Benefit-cost ratio (NY-LA)} = \frac{\$15\text{B/yr} \times 20 \text{ yrs (NPV)}}{{\$253\text{B}} + \$0.04\text{B}} \approx 1.2\times
\end{equation}

The relay network alone, at \$40 million for the NY-LA corridor, has a benefit-cost ratio
that is essentially unbounded in the first-year analysis: \$15 billion in annual benefits
against \$40 million in infrastructure cost implies a payback period under three weeks.
The capital constraint is not the relay stations. It is the managed lane program. The relay
stations are the unlock mechanism --- nearly free, foundational, and deployable now.

This is the core policy insight. The \$253 billion I2.0 program is the primary investment.
Its economic return depends partly on the relay network making 48-hour SLA achievable and
commercially reliable. The relay network is Layer 0 of I2.0 --- buildable first, at
minimal cost, activating economic returns before the main program is complete.
